\documentclass[conference]{IEEEtran}
\input epsf
\usepackage{graphicx}
\usepackage{amsmath,amssymb}
\usepackage{booktabs}
\usepackage[brazil]{babel}

\hyphenation{ba-lan-cea-dor re-qui-si-coes sim-u-la-cao}

\begin{document}

\title{\LARGE Balanceador de Carga para Sistema Distribuído:\\
Modelagem Analítica e Verificação por Simulação}

\author{\authorblockN{[Nome do Integrante 1] e [Nome do Integrante 2]}
\authorblockA{MC714 -- Sistemas Distribuídos\\
Instituto de Computação -- UNICAMP\\
Campinas, SP, Brasil}}

\maketitle

\begin{abstract}
Este trabalho implementa e analisa um balanceador de carga para um sistema distribuído composto por três servidores homogêneos M/M/1. O simulador de eventos discretos (SimPy) suporta as políticas Aleatória, Round-Robin e Fila Mais Curta (JSQ). Desenvolve-se modelagem analítica baseada no teorema de decomposição de Poisson e na fila M/M/1, comparando-se os valores teóricos com a média de dez execuções independentes da simulação.
\end{abstract}

\begin{keywords}
Balanceamento de carga, filas M/M/1, simulação de eventos discretos, teoria de filas.
\end{keywords}

%==============================================================================
\section{Resumo do Projeto}
%==============================================================================

O objetivo deste projeto é implementar, modelar analiticamente e validar por simulação um \textbf{balanceador de carga} que distribui requisições entre três servidores homogêneos em um sistema distribuído. O trabalho faz parte da disciplina MC714 (Sistemas Distribuídos) e aborda políticas clássicas de escalonamento, métricas de desempenho (vazão~$X$, tempo médio de resposta~$E[R]$, número médio no sistema~$E[N]$ e utilização~$U_i$) e a comparação rigorosa entre teoria e experimento.

A \textbf{arquitetura} adotada separa três componentes principais, conforme ilustrado conceitualmente na Fig.~\ref{fig:arquitetura}:

\begin{itemize}
    \item \textbf{Gerador de tráfego:} processo de chegadas Poisson($\lambda$), com tempos entre chegadas exponenciais de média~$1/\lambda$.
    \item \textbf{Balanceador:} componente instantâneo que, a cada requisição, seleciona um dos três servidores segundo a política configurada (Aleatória, Round-Robin ou JSQ).
    \item \textbf{Camada de servidores:} três filas M/M/1 homogêneas em FCFS, cada uma com taxa de serviço~$\mu = 1{,}0$~req/u.t.
\end{itemize}

O simulador foi implementado em Python com SimPy. Cada experimento principal executa $5000$~u.t.\ de simulação (com warm-up de $500$~u.t.), repetido $10$~vezes para obtenção de médias amostrais. Paralelamente, deduz-se um \textbf{modelo analítico} M/M/1 por servidor sob roteamento aleatório equiprovável ($p_i = 1/3$), permitindo comparar quantitativamente as previsões teóricas com os resultados simulados da política Aleatória e avaliar o ganho relativo das políticas informadas.

\begin{figure}[!t]
\centering
\fbox{\parbox{0.9\columnwidth}{\centering\vspace{0.3cm}
\textbf{Cliente} $\xrightarrow{\;\lambda\;}$ \textbf{Balanceador} $\xrightarrow{\;1/3\;}$ \textbf{Servidor 1}\\
\hspace{2.8cm}$\xrightarrow{\;1/3\;}$ \textbf{Servidor 2}\\
\hspace{2.8cm}$\xrightarrow{\;1/3\;}$ \textbf{Servidor 3}
\vspace{0.3cm}}}
\caption{Arquitetura do sistema: chegadas Poisson($\lambda$), balanceador com três políticas e três servidores M/M/1 homogêneos.}
\label{fig:arquitetura}
\end{figure}

%==============================================================================
\section{Arquitetura do Balanceador e Políticas Adotadas}
%==============================================================================

O balanceador é implementado como um módulo \textit{stateless} de decisão instantânea: a cada chegada, invoca-se uma função de política que recebe a lista de servidores e retorna o destino da requisição. Não há fila no balanceador nem custo de roteamento; a requisição é enfileirada imediatamente no servidor escolhido. A política ativa é selecionada por parâmetro de configuração, permitindo comparar as três estratégias sob condições idênticas de tráfego.

\subsection{Políticas Implementadas}

\textbf{Aleatória (Random):} a cada requisição, sorteia-se uniformemente um entre os três servidores ($p_i = 1/3$). Não mantém estado entre decisões. Esta política serve de referência tanto para a simulação quanto para o modelo analítico M/M/1.

\textbf{Round-Robin (RR):} mantém um contador cíclico que atribui requisições sequencialmente aos servidores $1 \to 2 \to 3 \to 1 \to \cdots$. Garante distribuição perfeitamente equitativa a cada ciclo de três requisições, eliminando desvios estocásticos da política aleatória.

\textbf{Fila Mais Curta (JSQ):} consulta a carga instantânea de cada servidor---número de requisições na fila mais a requisição em serviço---e encaminha ao servidor de menor carga. Empates são resolvidos aleatoriamente. Monitora ativamente o estado do sistema, nivelando filas em tempo real.

A Tabela~\ref{tab:politicas} resume as características comparativas.

\begin{table}[!t]
\centering
\caption{Comparação das políticas de balanceamento.}
\label{tab:politicas}
\small
\begin{tabular}{@{}lccc@{}}
\toprule
\textbf{Política} & \textbf{Estado} & \textbf{Informação} & \textbf{Ref.\ analítica} \\
\midrule
Aleatória   & Não  & Nenhuma       & Sim ($p_i\!=\!1/3$) \\
Round-Robin & Sim  & Índice ciclo  & Não \\
JSQ         & Não  & Carga filas   & Não \\
\bottomrule
\end{tabular}
\end{table}

\subsection{Implementação}

A arquitetura de software separa responsabilidades em módulos: \texttt{policies.py} define as três funções de roteamento; \texttt{simulator.py} implementa o loop de eventos discretos (SimPy) com a classe \texttt{Server} (fila FCFS, serviço exponencial); e \texttt{metrics.py} registra $N(t)$, tempos de resposta e utilização por servidor no intervalo $[500, 5000]$~u.t. O balanceador invoca a política no instante de cada chegada:

\begin{verbatim}
server = policy_fn(servers, rng, policy_state)
server.enqueue(request)
\end{verbatim}

Requisições rejeitadas por buffer finito (experimento bônus) são contabilizadas separadamente; nos experimentos principais, as filas são ilimitadas.

%==============================================================================
\section{Servidores e Parametrização do Tráfego}
%==============================================================================

\subsection{Modelo dos Servidores}

Cada um dos três servidores é modelado como uma fila \textbf{M/M/1} homogênea: um único canal de atendimento com disciplina FCFS (\textit{First-Come, First-Served}). O tempo de serviço segue distribuição exponencial com taxa $\mu = 1{,}0$~req/u.t., ou seja, $E[S] = 1$~u.t. A carga de um servidor é definida como o número de requisições aguardando na fila somado a $1$ se houver requisição em processamento. Nos experimentos principais, o buffer é \textbf{ilimitado}; requisições nunca são descartadas.

O estado de cada servidor (comprimento da fila e ocupação) é monitorado continuamente durante toda a simulação. A utilização $U_i$ de cada servidor~$i$ é estimada como a fração do intervalo de medição em que o servidor permanece ocupado.

\subsection{Parametrização do Tráfego}

As chegadas ao balanceador seguem um \textbf{processo de Poisson} de taxa~$\lambda$: os tempos entre chegadas consecutivas são independentes e exponencialmente distribuídos com média~$1/\lambda$. Este modelo estacionário permite confrontar diretamente os resultados simulados com as fórmulas analíticas M/M/1.

Os experimentos principais variam $\lambda$ em cinco níveis que cobrem regimes de baixa a alta carga ($\rho_i = \lambda/3 < 1$), além de um experimento de \textbf{instabilidade} com $\lambda = 3{,}3 > 3\mu$. Para cada combinação de $\lambda$ e política, executam-se \textbf{10 réplicas} independentes com sementes distintas; reportam-se as médias amostrais. A Tabela~\ref{tab:parametros} consolida os parâmetros.

\begin{table}[!t]
\centering
\caption{Parâmetros dos experimentos principais.}
\label{tab:parametros}
\small
\begin{tabular}{@{}ll@{}}
\toprule
\textbf{Parâmetro} & \textbf{Valor} \\
\midrule
Servidores ($n$)            & 3 homogêneos \\
Taxa de serviço ($\mu$)       & 1,0 req/u.t. \\
Taxas de chegada ($\lambda$)  & 0,6; 1,2; 1,8; 2,4; 2,7 \\
$\lambda$ instável            & 3,3 \\
Duração da simulação          & 5000 u.t. \\
Período de warm-up            & 500 u.t. \\
Intervalo de medição          & $[500, 5000]$ u.t. \\
Réplicas por configuração     & 10 \\
Configurações totais          & $5 \times 3 = 15$ \\
\bottomrule
\end{tabular}
\end{table}

\subsection{Métricas Coletadas}

No intervalo pós warm-up, calculam-se: \textbf{vazão}~$X$ (requisições completadas por u.t.), \textbf{tempo médio de resposta}~$E[R]$ (intervalo entre chegada e conclusão), \textbf{número médio no sistema}~$E[N]$ (via integração de $N(t)$) e \textbf{utilização média}~$U_i$ por servidor. Intervalos de confiança de 95\% ($t_{0{,}975,9}$) são computados sobre as dez réplicas, embora o relatório compare primariamente as médias amostrais com os valores analíticos.

%==============================================================================
\section{Modelagem Analítica}
%==============================================================================

\subsection{Item (a): Decomposição de Poisson (\textit{Splitting})}

Seja $N(t)$ um processo de Poisson($\lambda$) e $N_i(t)$ o número de requisições roteadas ao servidor~$i$ sob política aleatória com $p_i = 1/3$. Dado $N(t)=n$, a alocação $(n_1,n_2,n_3)$ segue uma multinomial; combinando com a PMF de Poisson e cancelando $n!$, obtém-se:
\begin{equation}
P(N_1\!=\!n_1, N_2\!=\!n_2, N_3\!=\!n_3) = \prod_{i=1}^{3} \frac{e^{-\lambda p_i t}(\lambda p_i t)^{n_i}}{n_i!}
\end{equation}
Logo, os processos $N_i(t)$ são \textbf{independentes} e cada um é Poisson($\lambda_i$), com $\lambda_i = \lambda/3$.

\subsection{Item (b): Fila M/M/1 por Servidor}

Com $\rho_i = \lambda_i/\mu = \lambda/(3\mu)$ e $\mu=1$, deduzem-se as métricas estacionárias por servidor e do sistema:
\begin{align}
U_i &= \rho_i, \quad E[N_i] = \frac{\rho_i}{1-\rho_i}, \quad
E[R] = \frac{1}{\mu - \lambda_i} = \frac{1}{\mu - \lambda/3}
\end{align}
Além disso, $E[T_Q] = E[R] - 1/\mu$, $X = \lambda$ e $E[N] = 3\,E[N_i]$. O sistema é \textbf{estável} se $\lambda < 3\mu$.

\subsection{Item (c): Conservação de Trabalho}

Sob estabilidade e filas ilimitadas, nenhuma requisição é perdida e servidores ociosos atendem imediatamente. Logo, $X = \lambda$ e $U_i = \lambda/3$ valem para \textbf{todas} as políticas, embora $E[R]$ e $E[N]$ dependam do roteamento.

\subsection{Item (d): Comparação de Políticas}

O modelo M/M/1 com splitting descreve apenas a política \textbf{Aleatória}. Políticas informadas reduzem desbalanceamento e satisfazem:
\begin{equation}
E[R]_{\text{JSQ}} \le E[R]_{\text{RR}} \le E[R]_{\text{Random}}
\end{equation}
Round-Robin elimina variações estocásticas locais; JSQ monitora filas e minimiza ociosidade.

\subsection{Item (e): Lei de Little}

Para qualquer sistema estável e conservativo, $E[N] = X \cdot E[R]$, independentemente da política. A Tabela~\ref{tab:little} (Apêndice) confirma erro relativo $< 0{,}05\%$ nas 15 configurações.

\subsection{Item (f): Caso Instável ($\lambda = 3{,}3$)}

Como $\lambda > 3\mu$, tem-se $\rho = 1{,}10 > 1$ e as fórmulas M/M/1 estacionárias divergem ($E[N] \to \infty$). A \textbf{aproximação fluida} prevê acúmulo linear:
\begin{equation}
N(t) \approx N(0) + (\lambda - 3\mu)\,t = N(0) + 0{,}3\,t
\end{equation}
A simulação confirma $X \approx 3{,}0$~req/u.t.\ (gargalo), $U_i \approx 1{,}0$ e $E[N] \approx 866$ na política aleatória, com crescimento linear de $N(t)$ conforme a Fig.~\ref{fig:unstable}.

%==============================================================================
\section{Resultados e Discussão}
%==============================================================================

\subsection{Validação Analítico $\times$ Simulação}

A Tabela~\ref{tab:comparacao} compara o modelo M/M/1 (política Aleatória) com a média de dez réplicas. Os desvios em $E[R]$, $X$ e $E[N]$ permanecem abaixo de 3\%, confirmando a adequação do modelo para roteamento equiprovável. Pequenas diferenças decorrem de variância amostral e do período finito de medição.

\begin{table}[!t]
\centering
\caption{Comparação analítico $\times$ simulação (política Aleatória, $n=10$).}
\label{tab:comparacao}
\small
\setlength{\tabcolsep}{3pt}
\begin{tabular}{@{}rcccc@{}}
\toprule
$\lambda$ & $E[R]_{\text{ana}}$ & $E[R]_{\text{sim}}$ & $E[N]_{\text{ana}}$ & $E[N]_{\text{sim}}$ \\
\midrule
0,6 & 1,25 & 1,25 & 0,75 & 0,75 \\
1,2 & 1,67 & 1,68 & 2,00 & 2,01 \\
1,8 & 2,50 & 2,50 & 4,50 & 4,53 \\
2,4 & 5,00 & 4,91 & 12,00 & 11,82 \\
2,7 & 10,00 & 10,30 & 27,00 & 27,84 \\
\bottomrule
\end{tabular}
\end{table}

\subsection{Desempenho das Políticas}

A Fig.~\ref{fig:er} mostra $E[R]$ em função de~$\lambda$. JSQ apresenta melhor desempenho, seguida de Round-Robin; a curva analítica coincide com a simulação aleatória. Sob $\lambda=2{,}7$, JSQ reduz $E[R]$ em $\approx$60\% em relação à política aleatória. Em todas as configurações, $X \approx \lambda$ e $U_i \approx \lambda/3$, validando a conservação de trabalho.

\begin{figure}[!t]
\centering
\includegraphics[width=\columnwidth]{results/er_vs_lambda.png}
\caption{$E[R]$ analítico (Aleatória) e simulado para as três políticas.}
\label{fig:er}
\end{figure}

\begin{figure}[!t]
\centering
\includegraphics[width=\columnwidth]{results/N_t_lambda33.png}
\caption{Evolução de $N(t)$ para $\lambda=3{,}3$ e reta fluida $N(0)+0{,}3t$.}
\label{fig:unstable}
\end{figure}

%==============================================================================
\section{Conclusão}
%==============================================================================

Implementou-se um simulador de balanceador de carga com três políticas clássicas sobre servidores M/M/1 homogêneos. A modelagem analítica baseada no teorema de splitting e na fila M/M/1 prediz corretamente o comportamento da política Aleatória, com concordância estatística nas métricas $X$, $E[R]$ e $E[N]$. Políticas informadas (Round-Robin e JSQ) reduzem significativamente o tempo de resposta sem alterar vazão nem utilização média. A Lei de Little foi verificada empiricamente e o regime instável ($\lambda=3{,}3$) foi coerente com a aproximação fluida. Trabalhos futuros incluem extensão a servidores heterogêneos e buffers finitos (implementados como bônus).

%==============================================================================
\appendices
\section{Tabelas Complementares}
%==============================================================================

\begin{table*}[!t]
\centering
\caption{Modelo analítico M/M/1 por servidor ($\mu=1{,}0$, política Aleatória).}
\label{tab:analitico}
\small
\begin{tabular}{@{}cccccc@{}}
\toprule
$\lambda$ & $\rho_i$ & $U_i$ & $E[N_i]$ & $E[R]$ & $X$ \\
\midrule
0,6 & 0,20 & 0,20 & 0,25 & 1,25 & 0,60 \\
1,2 & 0,40 & 0,40 & 0,67 & 1,67 & 1,20 \\
1,8 & 0,60 & 0,60 & 1,50 & 2,50 & 1,80 \\
2,4 & 0,80 & 0,80 & 4,00 & 5,00 & 2,40 \\
2,7 & 0,90 & 0,90 & 9,00 & 10,00 & 2,70 \\
\bottomrule
\end{tabular}
\end{table*}

\begin{table*}[!t]
\centering
\caption{Resultados simulados completos (média de 10 réplicas).}
\label{tab:simulacao}
\small
\setlength{\tabcolsep}{4pt}
\begin{tabular}{@{}lrrrrrr@{}}
\toprule
Política & $\lambda$ & $X$ & $E[R]$ & $E[N]$ & $U_{\text{avg}}$ & $E[R]_{\text{ana}}$ \\
\midrule
Aleatória      & 0,6 & 0,599 & 1,253 & 0,751 & 0,200 & 1,250 \\
Aleatória      & 1,2 & 1,196 & 1,682 & 2,012 & 0,400 & 1,667 \\
Aleatória      & 1,8 & 1,805 & 2,505 & 4,527 & 0,601 & 2,500 \\
Aleatória      & 2,4 & 2,409 & 4,905 & 11,817 & 0,800 & 5,000 \\
Aleatória      & 2,7 & 2,703 & 10,300 & 27,840 & 0,901 & 10,000 \\
Round-Robin    & 0,6 & 0,603 & 1,057 & 0,637 & 0,200 & 1,250 \\
Round-Robin    & 1,2 & 1,203 & 1,293 & 1,555 & 0,402 & 1,667 \\
Round-Robin    & 1,8 & 1,803 & 1,829 & 3,298 & 0,603 & 2,500 \\
Round-Robin    & 2,4 & 2,412 & 3,495 & 8,434 & 0,803 & 5,000 \\
Round-Robin    & 2,7 & 2,694 & 6,461 & 17,408 & 0,897 & 10,000 \\
JSQ            & 0,6 & 0,605 & 1,031 & 0,624 & 0,203 & 1,250 \\
JSQ            & 1,2 & 1,198 & 1,139 & 1,365 & 0,400 & 1,667 \\
JSQ            & 1,8 & 1,803 & 1,431 & 2,580 & 0,603 & 2,500 \\
JSQ            & 2,4 & 2,404 & 2,286 & 5,497 & 0,801 & 5,000 \\
JSQ            & 2,7 & 2,704 & 4,138 & 11,202 & 0,902 & 10,000 \\
\bottomrule
\end{tabular}
\end{table*}

\begin{table*}[!t]
\centering
\caption{Verificação da Lei de Little ($E[N] = X \cdot E[R]$).}
\label{tab:little}
\small
\setlength{\tabcolsep}{4pt}
\begin{tabular}{@{}lrrrrr@{}}
\toprule
Política & $\lambda$ & $X \cdot E[R]$ & $E[N]$ medido & Erro (\%) \\
\midrule
Aleatória   & 0,6 & 0,751 & 0,751 & 0,006 \\
Aleatória   & 1,2 & 2,012 & 2,012 & 0,001 \\
Aleatória   & 1,8 & 4,524 & 4,527 & 0,071 \\
Aleatória   & 2,4 & 11,816 & 11,817 & 0,012 \\
Aleatória   & 2,7 & 27,863 & 27,840 & 0,082 \\
Round-Robin & 0,6 & 0,637 & 0,637 & 0,022 \\
Round-Robin & 1,2 & 1,555 & 1,555 & 0,010 \\
Round-Robin & 1,8 & 3,299 & 3,298 & 0,014 \\
Round-Robin & 2,4 & 8,429 & 8,434 & 0,054 \\
Round-Robin & 2,7 & 17,408 & 17,408 & 0,000 \\
JSQ         & 0,6 & 0,624 & 0,624 & 0,002 \\
JSQ         & 1,2 & 1,365 & 1,365 & 0,013 \\
JSQ         & 1,8 & 2,580 & 2,580 & 0,011 \\
JSQ         & 2,4 & 5,498 & 5,497 & 0,010 \\
JSQ         & 2,7 & 11,196 & 11,202 & 0,053 \\
\bottomrule
\end{tabular}
\end{table*}

\end{document}
